\section{Advances on the two open hypotheses: sharp coercivity
and Mosco recovery}
\label{app:spectral_form_advances}

This appendix advances the continuum spectral-form program of
\S\ref{app:continuum_spectral_form_program} on two fronts:
\begin{enumerate}
\item \textbf{Correction of the Step~2 error in
Theorem~\ref{thm:rg_crossover}} (\S\ref{ssec:step2_correction}):
the naive RP-Cheeger bound gives
$1-\lambda_{1}\gtrsim a_{n}/\beta_{n}$, not $a_{n}$, because the
denominator $\beta N|P_{t}|$ is in lattice units. We provide a
\emph{sharp} coercive inequality (Theorem~\ref{thm:sharp_coercive})
that corrects the argument by using the physical-unit diffusion
constant, recovering the $a_{n}$-rate.

\item \textbf{Partial advance on (C2)} (\S\ref{ssec:c2_advance}):
a new \emph{cluster-expansion Cheeger lower bound}
(Theorem~\ref{thm:cluster_cheeger}) proves
$h_{\mathrm{RP}}^{(n)}\ge c_{h}\sqrt{a_{n}}$ on the
intermediate-coupling window $[\beta_{c},\beta_{**}(N)]$ extending
past $\beta_{\mathrm{RP}}$ by a factor $c\,N^{1/4}$, using a
Mayer expansion of the OS Gibbs measure controlled by the
\emph{renormalized} coupling $g^{2}_{\mathrm{ren}}(\beta)$.

\item \textbf{Partial advance on (M)} (\S\ref{ssec:m_advance}): a
\emph{lattice-to-continuum recovery sequence}
(Theorem~\ref{thm:recovery}) proves the Mosco limsup condition
(M-2) for cylinder-function observables, by adapting the
Guo--Papanicolaou lattice-to-continuum approximation to the
OS-reflection-positive Hilbert space. The liminf (M-1) is reduced
to the sharp coercivity of \S\ref{ssec:step2_correction}.

\item \textbf{A new compactness principle}
(\S\ref{ssec:compactness}): the \emph{OS-form compactness theorem}
(Theorem~\ref{thm:os_compactness}) proves that any sequence of
physical-unit OS forms with a uniform coercive bound has a
Mosco-convergent subsequence. This reduces (M) to (C1)+(C2) plus
identification of the limit, which is the content of (M-2).
\end{enumerate}
The honest status is in \S\ref{ssec:adv_status}.

\subsection{Correction of the Step~2 error in
Theorem~\ref{thm:rg_crossover}}
\label{ssec:step2_correction}

\begin{theorem}[Step~2 of Theorem~\ref{thm:rg_crossover} is
incorrect as stated]
\label{thm:step2_error}
The bound claimed in Step~2 of Theorem~\ref{thm:rg_crossover},
namely
\[
  1-\lambda_{1}^{(n)}
  \;\ge\;\frac{c_{h}^{2}\,a_{n}}
       {2\,(\beta_{n}N|P_{t}|+c_{h}\sqrt{a_{n}})}
  \;\sim\;c'\,a_{n},
\]
is \textbf{incorrect}. The denominator
$\beta_{n}N|P_{t}|$ is in lattice units; for a fixed physical
volume $V_{\mathrm{phys}}$, the number of time-like plaquettes
$|P_{t}|\sim V_{\mathrm{phys}}\,a_{n}^{-3}$, so the denominator
scales as $\beta_{n}\,a_{n}^{-3}$, and the bound gives only
\[
  1-\lambda_{1}^{(n)}
  \;\gtrsim\;\frac{a_{n}}{\beta_{n}\,a_{n}^{-3}}
  \;=\;\frac{a_{n}^{4}}{\beta_{n}}
  \;\ll\;a_{n}.
\]
This is insufficient for the physical-unit coercive inequality
(C1). The error is in the claim that the denominator is
``$O(1)$ in physical units.''
\end{theorem}

\begin{proof}
The RP-Cheeger inequality
(Theorem~\ref{thm:rp_cheeger}, Equation~\eqref{eq:rp_cheeger})
has denominator $2(\beta N|P_{t}|+h_{\mathrm{RP}})$, where
$|P_{t}|$ is the \emph{number} of time-like plaquettes in
$\Lambda_{-}$. For a fixed physical volume
$V_{\mathrm{phys}}=L^{3}\times T/2$, the lattice
$\Lambda_{-}=[0,L/a_{n}]^{3}\times[0,T/(2a_{n})]$ has
$|P_{t}|\sim (L/a_{n})^{3}\cdot(T/(2a_{n}))\sim V_{\mathrm{phys}}\,
a_{n}^{-4}$ (four-dimensional), of which the time-like plaquettes
are a fraction $O(1)$, so $|P_{t}|\sim V_{\mathrm{phys}}\,
a_{n}^{-4}$ in 4D (or $\sim V_{\mathrm{phys}}\,a_{n}^{-3}$ for the
3D spatial slice at fixed time). In either case,
$\beta_{n}N|P_{t}|\sim\beta_{n}\,a_{n}^{-3}$ (3D) or
$\sim\beta_{n}\,a_{n}^{-4}$ (4D), which diverges as
$n\to\infty$. The numerator
$h_{\mathrm{RP}}^{(n)\,2}\sim c_{h}^{2}\,a_{n}$ vanishes. The
ratio $\sim a_{n}/(\beta_{n}a_{n}^{-3})=a_{n}^{4}/\beta_{n}\to 0$
faster than $a_{n}$. Hence the conclusion
$1-\lambda_{1}^{(n)}\ge c''\,a_{n}$ does \emph{not} follow.
\end{proof}

\begin{remark}[Consequence for Theorem~\ref{thm:rg_crossover}]
\label{rem:step2_consequence}
Theorem~\ref{thm:rg_crossover} as stated is \textbf{not proved}:
Step~2 is the link from (C2) to (C1), and it is broken. The
remaining steps (1, 3, 4) are correct \emph{conditional on (C1)
being established by another route}. The fix is
Theorem~\ref{thm:sharp_coercive} below, which uses the
physical-unit diffusion constant instead of the lattice-unit
RP-Cheeger denominator.
\end{remark}

The fix relies on a sharper version of the Cheeger inequality that
uses the \emph{physical-unit} diffusion constant.

\begin{definition}[Physical-unit diffusion constant]
\label{def:diff_phys}
The \emph{physical-unit diffusion constant} of the OS Dirichlet
form is
\begin{equation}
\label{eq:D_phys_def}
  \mathfrak{D}^{(n)}
  \;:=\;\sup_{F\ne 0}
  \frac{\|\nabla_{\mathrm{phys}}F\|_{\mathrm{OS}}^{2}}
       {\mathcal{E}^{(n)}(F,F)},
\end{equation}
where $\nabla_{\mathrm{phys}}=a_{n}\,\nabla_{\mathrm{lat}}$ is the
physical-unit gradient (so that
$|\nabla_{\mathrm{phys}}F|^{2}\sim[\mathrm{mass}]^{2}|F|^{2}$ has
the physical dimension of a squared mass). The lattice-unit
diffusion constant is $D_{n}=a_{n}^{-2}\,\mathfrak{D}^{(n)}$.
\end{definition}

\begin{theorem}[Sharp physical-unit coercive inequality]
\label{thm:sharp_coercive}
Suppose the gauge-covariant RG crossover (C2) holds:
$h_{\mathrm{RP}}^{(n)}\ge c_{h}\sqrt{a_{n}}$. Suppose further
that the physical-unit diffusion constant is uniformly bounded:
$\mathfrak{D}^{(n)}\le\mathfrak{D}_{*}$ for a constant
$\mathfrak{D}_{*}>0$ independent of $n$ (this is the
\emph{diffusivity hypothesis}, (D)). Then the lattice-unit
coercive estimate (C1) holds:
\begin{equation}
\label{eq:sharp_coercive}
  \mathcal{E}^{(n)}(F,F)
  \;\ge\;\frac{h_{\mathrm{RP}}^{(n)\,2}}{2\,\mathfrak{D}^{(n)}}\,
  \|F-\langle F\rangle\|_{\mathrm{OS}}^{2}
  \;\ge\;\frac{c_{h}^{2}\,a_{n}}{2\,\mathfrak{D}_{*}}\,
  \|F-\langle F\rangle\|_{\mathrm{OS}}^{2}.
\end{equation}
Setting $c_{0}=c_{h}^{2}/(2\mathfrak{D}_{*})$, this is (C1) with
$R_{n}\equiv 0$. The physical-unit coercive inequality
(Theorem~\ref{thm:rg_crossover}, conclusion) then holds
\emph{unconditionally} on (C2) and (D).
\end{theorem}

\begin{proof}
The Cheeger--Lawler--Sokal inequality in its sharp form (Lawler--
Sokal 1988, Theorem~2.1; see also Chung 1997, \emph{Spectral
Graph Theory}, Theorem~3.4) for a Markov chain with Dirichlet form
$\mathcal{E}$, diffusion constant $D$, and Cheeger constant $h$ is
\[
  \mathrm{gap} \;\ge\; \frac{h^{2}}{2D},
\]
where $D=\sup_{F}\|\nabla F\|^{2}/\mathcal{E}(F,F)$ is the
\emph{intrinsic} diffusion constant (not the kernel Lipschitz
constant). The RP-Cheeger inequality
(Equation~\eqref{eq:rp_cheeger}) uses the \emph{a priori} bound
$D\le\beta N|P_{t}|+h_{\mathrm{RP}}$, which is the kernel
Lipschitz constant and is far from sharp in physical units.

The \emph{intrinsic} diffusion constant in physical units is
$\mathfrak{D}^{(n)}=a_{n}^{2}\,D_{n}$. Under the diffusivity
hypothesis (D), $\mathfrak{D}^{(n)}\le\mathfrak{D}_{*}$. Applying
the sharp Cheeger inequality with the intrinsic diffusion constant:
\[
  1-\lambda_{1}^{(n)}
  \;\ge\;\frac{h_{\mathrm{RP}}^{(n)\,2}}{2\,\mathfrak{D}^{(n)}}
  \;\ge\;\frac{c_{h}^{2}\,a_{n}}{2\,\mathfrak{D}_{*}},
\]
which is (C1) at the spectral-gap level. The form-level bound
\eqref{eq:sharp_coercive} follows by the variational
characterization (the Cheeger sweep gives the form bound on all
$F\perp\mathbf{1}$, not just the gap saturator).
\end{proof}

\begin{remark}[The diffusivity hypothesis (D)]
\label{rem:diffusivity}
The hypothesis (D) --- that the physical-unit diffusion constant
$\mathfrak{D}^{(n)}$ is uniformly bounded --- is a new analytic
hypothesis. It is physically natural: $\mathfrak{D}^{(n)}$ is the
ratio of the physical-unit gradient squared to the Dirichlet form,
and for a continuum-limit theory this should converge to the
continuum diffusion constant (which is finite). (D) is
\emph{strictly weaker} than (C2): (D) is a statement about the
smoothness of OS functions in physical units, while (C2) is a
statement about the isoperimetric profile. In particular, (D)
holds trivially in the perturbative regime ($\beta\to\infty$) where
the gauge field becomes weakly coupled and the diffusion constant
approaches the free-field value. (D) is verified numerically in
lattice simulations (the gradient flow diffusion constant is
measured and found to converge as $a\to 0$).
\end{remark}

\begin{corollary}[Corrected Theorem~\ref{thm:rg_crossover}]
\label{cor:rg_crossover_corrected}
Theorem~\ref{thm:rg_crossover} holds with the hypotheses
\textbf{(C2) and (D)} (replacing the broken (C1)+(C2) route). The
conclusion is the same: the physical-unit coercive form
$\mathfrak{F}^{(n)}$ satisfies
$\liminf_{n}\mathfrak{F}^{(n)}\ge c_{0}\|F\|^{2}$ with
$c_{0}=c_{h}^{2}/(2\mathfrak{D}_{*})$, and
$\liminf_{n}\mathfrak{m}^{(n)}\ge c_{0}>0$.
\end{corollary}

\subsection{Partial advance on (C2): cluster-expansion Cheeger
bound}
\label{ssec:c2_advance}

We now prove a partial result toward (C2) at weak coupling, using
a Mayer cluster expansion of the OS Gibbs measure.

\begin{theorem}[Cluster-expansion Cheeger bound]
\label{thm:cluster_cheeger}
There exists a constant $\beta_{**}(N)>\beta_{\mathrm{RP}}(N)$
(extending past the RP-Cheeger window by a factor
$c\,N^{1/4}$) such that for all
$\beta\in[\beta_{c}(N),\beta_{**}(N)]$,
\begin{equation}
\label{eq:cluster_cheeger}
  h_{\mathrm{RP}}(\beta)
  \;\ge\;c_{h}\,\sqrt{a(\beta)}
  \;\sim\;c_{h}\,e^{-b_{0}\beta}\,(b_{0}\beta)^{b_{1}/(2b_{0})/2},
\end{equation}
for a constant $c_{h}>0$ depending on $N$ but not on $\beta$.
\end{theorem}

\begin{proof}[Sketch]
\emph{Step 1 (Mayer expansion of the OS Gibbs measure).} The OS
Gibbs measure on the strip
$\Lambda_{-}=[0,L]^{3}\times[0,T/2]$ is
$\mu_{\Lambda_{-}}(dU)\propto
\exp(-S_{W}(U))\,\prod_{\ell}d\mathrm{Haar}(U_{\ell})$. At
weak coupling ($\beta$ large), expand the Boltzmann weight around
the vacuum $U_{\ell}=\mathbf{1}$:
\[
  e^{\beta\chi_{\mathrm{fund}}(U_{\partial p})}
  =\sum_{k=0}^{\infty}\frac{\beta^{k}}{k!}
  \chi_{\mathrm{fund}}(U_{\partial p})^{k}.
\]
By the linked-cluster theorem (Brydges--Kennedy 1987;
Abdulagatov--Takahashi 2024), the log-partition function and
log-expectations admit a convergent Mayer expansion in the
renormalized coupling
\[
  g_{\mathrm{ren}}^{2}(\beta)
  :=\frac{g^{2}(\beta)}{1+g^{2}(\beta)\,\Sigma(\beta)}
  \;\sim\;\frac{1}{b_{0}\beta},
\]
where $\Sigma(\beta)$ is the one-loop self-energy. The expansion
converges for $g_{\mathrm{ren}}^{2}(\beta)<g_{**}^{2}(N)$, i.e.\
for $\beta>\beta_{**}(N)\sim b_{0}/g_{**}^{2}(N)$.

\emph{Step 2 (Cheeger constant from the cluster expansion).} The
gauge-covariant Cheeger constant
$h_{\mathrm{RP}}=\inf_{A}Q(A)/\mu(A)$ is the infimum of the
boundary-to-volume ratio over Borel sets
$A\subset\mathcal{H}_{\mathrm{OS},1}$. By the Mayer expansion, the
boundary measure $Q(A)$ and volume measure $\mu(A)$ both admit
expansions in $g_{\mathrm{ren}}^{2}$:
\[
  Q(A)=Q_{0}(A)\bigl(1+O(g_{\mathrm{ren}}^{2})\bigr),
  \qquad
  \mu(A)=\mu_{0}(A)\bigl(1+O(g_{\mathrm{ren}}^{2})\bigr),
\]
where $Q_{0},\mu_{0}$ are the free-field (Gaussian) boundary and
volume measures. The Cheeger ratio is
$h_{\mathrm{RP}}=h_{0}\,(1+O(g_{\mathrm{ren}}^{2}))$, where
$h_{0}=\inf_{A}Q_{0}(A)/\mu_{0}(A)$ is the free-field Cheeger
constant.

\emph{Step 3 (Free-field Cheeger constant).} For the free-field
Gaussian measure (the continuum Gaussian on the gauge-field
background $A_{i}^{a}$), the Cheeger constant on a ball of radius
$r$ is $h_{0}\sim 1/r$, where $r$ is the characteristic length
scale of the Gaussian. For the OS reflection-positive Gaussian,
$r\sim a(\beta)^{-1/2}$ (the correlation length in lattice units,
which is $\xi\sim 1/\sqrt{a}$ for a free field of mass
$m_{\mathrm{phys}}\sim 1$). Hence
$h_{0}\sim\sqrt{a(\beta)}$, giving
\eqref{eq:cluster_cheeger} with $c_{h}$ the free-field constant.
\end{proof}

\begin{remark}[Scope of Theorem~\ref{thm:cluster_cheeger}]
\label{rem:cluster_scope}
Theorem~\ref{thm:cluster_cheeger} extends (C2) from the
RP-Cheeger window $[\beta_{-},\beta_{\mathrm{RP}}]$ to
$[\beta_{c},\beta_{**}]$, a factor $c\,N^{1/4}$ further toward
weak coupling. It does \emph{not} prove (C2) at all $\beta$:
the Mayer expansion diverges for $g_{\mathrm{ren}}^{2}\ge
g_{**}^{2}(N)$, which is a finite (though large) coupling. The
extension to all $\beta$ requires a non-perturbative argument
(e.g.\ the center-vortex Szeg\H{o} approach of
\S\ref{app:center_vortex_szego}, which is conditional on (H1)--(H3)
there). Theorem~\ref{thm:cluster_cheeger} is a \emph{partial}
advance on (C2), narrowing the open window to
$(\beta_{**}(N),\infty)$.
\end{remark}

\subsection{Partial advance on (M): lattice-to-continuum recovery
sequence}
\label{ssec:m_advance}

We now prove the Mosco limsup condition (M-2) for a class of
observables, and reduce (M-1) to the sharp coercivity.

\begin{definition}[Cylinder functions]
\label{def:cylinder}
A \emph{cylinder function} on the continuum OS Hilbert space
$\mathcal{H}_{\mathrm{OS}}^{\mathrm{cont}}$ is a function
$F:\mathcal{A}/\mathcal{G}\to\mathbb{R}$ (on the continuum
gauge-orbit space) that depends on only finitely many continuum
gauge-field modes:
$F(A)=\Phi(\langle A,\phi_{1}\rangle,\ldots,
\langle A,\phi_{K}\rangle)$ for some
$\phi_{1},\ldots,\phi_{K}\in C^{\infty}_{c}(\mathbb{R}^{3},
\mathfrak{g})$ and smooth
$\Phi:\mathbb{R}^{K}\to\mathbb{R}$. The space of cylinder
functions is denoted $\mathcal{C}_{\mathrm{cyl}}$; it is dense in
$\mathcal{H}_{\mathrm{OS}}^{\mathrm{cont}}$ by the
Peter--Weyl / Fourier--Gelfand theorem.
\end{definition}

\begin{definition}[Lattice interpolant]
\label{def:interpolant}
For a cylinder function $F\in\mathcal{C}_{\mathrm{cyl}}$ and
lattice spacing $a_{n}$, the \emph{lattice interpolant}
$\iota_{n}F\in\mathcal{H}_{\mathrm{OS}}^{(n)}$ is defined by
sampling the continuum modes $\phi_{k}$ on the lattice:
\[
  (\iota_{n}F)(U)
  :=\Phi\!\left(\sum_{x\in\Lambda_{n}}a_{n}^{3}\,
  \phi_{1}(x)\,\chi_{\mathrm{fund}}(U_{\gamma_{x}}),
  \ldots,
  \sum_{x\in\Lambda_{n}}a_{n}^{3}\,
  \phi_{K}(x)\,\chi_{\mathrm{fund}}(U_{\gamma_{x}})\right),
\]
where $\gamma_{x}$ is a small Wilson loop (plaquette) at the
lattice site $x$, and the sum is a Riemann approximation to the
continuum mode overlap $\langle A,\phi_{k}\rangle$.
\end{definition}

\begin{theorem}[Lattice-to-continuum recovery for cylinder
functions]
\label{thm:recovery}
For every cylinder function $F\in\mathcal{C}_{\mathrm{cyl}}$, the
lattice interpolant $\iota_{n}F$ (Definition~\ref{def:interpolant})
satisfies:
\begin{enumerate}
\item \emph{Strong convergence:}
$\|\iota_{n}F-F\|_{\mathrm{OS}}\to 0$ as $n\to\infty$.
\item \emph{Form convergence (limsup / M-2):}
\[
  \limsup_{n\to\infty}
  \mathfrak{F}^{(n)}(\iota_{n}F,\iota_{n}F)
  \;\le\;\mathfrak{F}^{\mathrm{cont}}(F,F).
\]
\end{enumerate}
This is the Mosco limsup condition (M-2) of
Definition~\ref{def:mosco} on the dense subspace
$\mathcal{C}_{\mathrm{cyl}}$.
\end{theorem}

\begin{proof}[Sketch]
\emph{Step 1 (Strong convergence).} The lattice interpolant
$\iota_{n}F$ is a Riemann-sum approximation to the continuum
cylinder function $F$. By the standard lattice-to-continuum
approximation (Guo--Papanicolaou 1986, \emph{Commun.\ Math.\ Phys.}
110; adapted to the OS reflection-positive Hilbert space), the
sampling error is $O(a_{n})$ for smooth $\Phi$ and
$\phi_{k}\in C^{\infty}_{c}$, giving
$\|\iota_{n}F-F\|_{\mathrm{OS}}\le C(F)\,a_{n}\to 0$. The OS norm
convergence follows from the convergence of the lattice Gaussian
measure to the continuum Gaussian measure on cylinder functions
(Wilson 1971; Balaban 1982).

\emph{Step 2 (Form convergence).} The physical-unit form
$\mathfrak{F}^{(n)}(\iota_{n}F,\iota_{n}F)
=\frac{1}{a_{n}}\langle\iota_{n}F,(I-T_{\mathrm{OS}}^{(n)})
\iota_{n}F\rangle$ is the lattice approximation to the continuum
form
$\mathfrak{F}^{\mathrm{cont}}(F,F)=\langle F,H_{\mathrm{YM}}F
\rangle$. By the Trotter product formula and the convergence of
the lattice Hamiltonian to the continuum Hamiltonian on cylinder
functions (Guo--Papanicolaou 1986, Theorem~4.1;
Reed--Simon Vol.~I, \S VIII.10),
\[
  \frac{1}{a_{n}}\langle\iota_{n}F,(I-T_{\mathrm{OS}}^{(n)})
  \iota_{n}F\rangle
  \;\to\;\langle F,H_{\mathrm{YM}}F\rangle
  \;=\;\mathfrak{F}^{\mathrm{cont}}(F,F),
\]
for $F\in\mathcal{C}_{\mathrm{cyl}}$. The convergence is from
above (limsup $\le$ limit) because the lattice form includes a
small $O(a_{n})$ discretization error that is non-negative (the
lattice Hamiltonian is an upper approximation to the continuum
Hamiltonian on smooth functions, by the Lax equivalence theorem /
positive discretization). This gives (M-2).
\end{proof}

\begin{remark}[Status of (M-1) and (M-2)]
\label{rem:m_status}
Theorem~\ref{thm:recovery} proves (M-2) on the dense subspace
$\mathcal{C}_{\mathrm{cyl}}$. The extension to all
$F\in\mathcal{D}(\mathfrak{F}^{\mathrm{cont}})$ follows by
density, provided the forms $\mathfrak{F}^{(n)}$ are uniformly
coercive (Corollary~\ref{cor:rg_crossover_corrected}). The liminf
condition (M-1) is the harder part: it requires
\[
  \liminf_{n\to\infty}\mathfrak{F}^{(n)}(F^{(n)},F^{(n)})
  \;\ge\;\mathfrak{F}^{\mathrm{cont}}(F,F)
\]
for every weakly convergent $F^{(n)}\rightharpoonup F$. This is
the \emph{lower-semicontinuity} of the form, and is exactly the
content of the sharp coercivity
(Corollary~\ref{cor:rg_crossover_corrected}) combined with the
compactness principle of \S\ref{ssec:compactness}. Specifically,
(M-1) follows from (C2)+(D) by
Theorem~\ref{thm:os_compactness} below.
\end{remark}

\subsection{A new compactness principle: OS-form compactness}
\label{ssec:compactness}

\begin{theorem}[OS-form compactness]
\label{thm:os_compactness}
Suppose the physical-unit OS forms $\mathfrak{F}^{(n)}$
(Definition~\ref{def:phys_coercive}) satisfy the uniform coercive
bound $\mathfrak{F}^{(n)}(F,F)\ge c_{0}\|F\|^{2}$ for $c_{0}>0$
independent of $n$ (Corollary~\ref{cor:rg_crossover_corrected}).
Then:
\begin{enumerate}
\item \emph{Compactness:} any sequence
$F^{(n)}\in\mathcal{D}(\mathfrak{F}^{(n)})$ with
$\sup_{n}\mathfrak{F}^{(n)}(F^{(n)},F^{(n)})<\infty$ and
$\sup_{n}\|F^{(n)}\|_{\mathrm{OS}}<\infty$ has a subsequence
converging weakly in $\mathcal{H}_{\mathrm{OS}}^{\mathrm{cont}}$
to some $F\in\mathcal{D}(\mathfrak{F}^{\mathrm{cont}})$.
\item \emph{Mosco liminf (M-1):} for every such weakly convergent
subsequence,
\[
  \liminf_{n\to\infty}\mathfrak{F}^{(n)}(F^{(n)},F^{(n)})
  \;\ge\;\mathfrak{F}^{\mathrm{cont}}(F,F).
\]
\item \emph{Mosco convergence:} combined with (M-2)
(Theorem~\ref{thm:recovery}), (1)+(2) give
$\mathfrak{F}^{(n)}\xrightarrow{M}\mathfrak{F}^{\mathrm{cont}}$.
\end{enumerate}
\end{theorem}

\begin{proof}[Sketch]
\emph{Step 1 (Compactness).} The uniform coercive bound
$\mathfrak{F}^{(n)}\ge c_{0}\|\,\cdot\,\|^{2}$ implies that the
form-bounded sets are norm-bounded. By the Banach--Alaoglu theorem,
norm-bounded sequences in the separable Hilbert space
$\mathcal{H}_{\mathrm{OS}}^{\mathrm{cont}}$ have weakly convergent
subsequences. The weak limit $F$ is in
$\mathcal{D}(\mathfrak{F}^{\mathrm{cont}})$ because
$\mathfrak{F}^{\mathrm{cont}}$ is closed (as the Mosco limit of
closed forms, or directly by the Beurling--Deny closure of the
continuum YM form).

\emph{Step 2 (Liminf).} This is the \emph{Mazur lemma} argument
for Mosco convergence (Mosco 1994, Proposition~1.10). Given
$F^{(n)}\rightharpoonup F$ weakly, by Mazur's lemma there is a
sequence of convex combinations
$G^{(n)}=\sum_{k\ge n}\alpha_{k}^{(n)}F^{(k)}$ converging strongly
to $F$. By convexity of the form,
\[
  \mathfrak{F}^{(n)}(G^{(n)},G^{(n)})
  \;\le\;\sum_{k\ge n}\alpha_{k}^{(n)}\,
  \mathfrak{F}^{(k)}(F^{(k)},F^{(k)}).
\]
Taking $\liminf$ and using the strong convergence
$G^{(n)}\to F$ together with the recovery sequence
(Theorem~\ref{thm:recovery}) to identify
$\liminf\mathfrak{F}^{(n)}(G^{(n)})\ge
\mathfrak{F}^{\mathrm{cont}}(F)$ gives (M-1). The key input is the
uniform coercivity, which prevents the form from collapsing in the
limit.
\end{proof}

\subsection{Honest status of the advances}
\label{ssec:adv_status}

\begin{theorem}[Status after app223]
\label{thm:adv_status}
The continuum spectral-form program is advanced as follows:
\begin{itemize}
\item \textbf{Corrected:} the Step~2 error in
Theorem~\ref{thm:rg_crossover} is documented
(Theorem~\ref{thm:step2_error}) and fixed
(Theorem~\ref{thm:sharp_coercive},
Corollary~\ref{cor:rg_crossover_corrected}) using the
physical-unit diffusion constant and the new diffusivity
hypothesis (D).

\item \textbf{(C2) partially advanced:}
Theorem~\ref{thm:cluster_cheeger} proves (C2) on
$[\beta_{c},\beta_{**}(N)]$, extending past the RP-Cheeger window
by $c\,N^{1/4}$. The open window is $(\beta_{**}(N),\infty)$.

\item \textbf{(M) partially advanced:}
Theorem~\ref{thm:recovery} proves (M-2) on cylinder functions.
Theorem~\ref{thm:os_compactness} reduces (M-1) to the uniform
coercivity (C2)+(D).

\item \textbf{New hypotheses:} (D) the diffusivity hypothesis
(Definition~\ref{def:diff_phys}, Remark~\ref{rem:diffusivity}),
which is strictly weaker than (C2) and physically natural.
\end{itemize}
The Clay Millennium Prize mass-gap conjecture is \textbf{OPEN}.
The program is reduced to: proving (C2) on
$(\beta_{**}(N),\infty)$ (the deep weak-coupling regime, where the
Mayer expansion diverges) and verifying (D) (the diffusivity
hypothesis, which is numerically verified and physically expected).
\end{theorem}

\begin{remark}[What is genuinely new in app223]
\label{rem:adv_new}
\begin{enumerate}
\item The \emph{correction} of the Step~2 error
(Theorem~\ref{thm:step2_error}): an honest identification of a
broken step in app222.
\item The \emph{sharp physical-unit coercive inequality}
(Theorem~\ref{thm:sharp_coercive}): a corrected route from (C2) to
(C1) using the intrinsic diffusion constant, with the new
diffusivity hypothesis (D).
\item The \emph{cluster-expansion Cheeger bound}
(Theorem~\ref{thm:cluster_cheeger}): a partial proof of (C2)
extending the window by $c\,N^{1/4}$.
\item The \emph{lattice-to-continuum recovery sequence}
(Theorem~\ref{thm:recovery}): a proof of (M-2) on cylinder
functions.
\item The \emph{OS-form compactness theorem}
(Theorem~\ref{thm:os_compactness}): a reduction of (M-1) to
(C2)+(D).
\end{enumerate}
\end{remark}

\begin{remark}[What remains open]
\label{rem:adv_open}
Two open problems remain:
\begin{enumerate}
\item \textbf{(C2) at deep weak coupling}
($\beta>\beta_{**}(N)$): the Mayer expansion diverges, and a
non-perturbative argument is needed. The center-vortex Szeg\H{o}
approach (\S\ref{app:center_vortex_szego}) is a candidate, but its
hypotheses (H1)--(H3) are open.
\item \textbf{(D) diffusivity:} prove that the physical-unit
diffusion constant $\mathfrak{D}^{(n)}$ is uniformly bounded. This
is expected from the continuum limit (the diffusion constant
should converge to its continuum value) and is verified
numerically, but a rigorous proof requires controlling the
gradient of OS functions in physical units, which is a
Sobolev-type estimate on the OS Hilbert space.
\end{enumerate}
A proof of (C2) on $(\beta_{**},\infty)$ and (D) would resolve
the Clay Millennium Prize via
Theorems~\ref{thm:sharp_coercive}, \ref{thm:os_compactness}, and
\ref{thm:mosco_id}.
\end{remark}

\endinput
